
In this paper, we paid equal attention to the substantive issue of how visual illusions covary across individuals and the methodological issue of how to document structure in this covariation.  The following conclusions may be drawn substantively about visual illusions:  

1. Visual illusions are strikingly large in effect and straightforward to document.  Data from each of our participants in each task, comprising at most 15 observations, often shows more evidence of an effect than typical cognitive psychology experiments comprised of hundreds of observations from scores of participants.  

2.  Individual differences are large.  The amount of variance across individuals is large relative to trial noise.  The ratio, $\gamma^2$ is around 1.145-to-1 (in variances), which is about 73 times larger than those in cognitive control tasks \parencite[]{Rouder.etal.2024}.  

3. Correlations across different illusions are relatively small.  This small size is not due to statistical attenuation but reflects the true state of things---illusions are at most weakly related.  

4. Nonetheless, there is at least one common factor, moreover, this factor accounts for 23.3\% of the total variance.

Our methodology for uncovering the structure of individual differences has a few elements that are not commonly stressed: 

1. We spend much effort assessing how much individual variation is recoverable.  The task plots in Figure~\ref{fig:taskPlots} show how different individuals are in each effect.  If these differences are hard to discern graphically, then it may be hard to document correlations among tasks.

2. Understanding the uncertainty in correlations is key to placing them in context and performing subsequent latent-variable modeling.  The plot in Figure \ref{fig:ciCorPlots} shows immediately that with the exclusion of version correlations, correlations are modest and fairly uniform, and moreover, there is a positive manifold across them.   The modesty of these correlations relative to uncertainty means that it is not possible to derive a many-factor solution.  One issue is that obtaining these uncertainties is not straightforward.  They come from a hierarchical model where the data model accounts for trial-to-trial variation and the model of individuals accounts for between-participant covariation.  Bayesian hierarchical models are perfect for this endeavor and there are a wide range of successful applications and tutorials for psychologists. 

3. Latent-variable modeling was difficult on several accounts.  The main problem was accounting for high version correlations for some illusions while accounting for lower version correlations in others.  We relied on the inspection of the correlations for guidance.

\subsection{Interpretation}

In the literature, there are seemingly two incompatible results.  The first is that individual differences in visual illusions are at best weakly correlated \parencite{Cretenoud.etal.2019} and the proportion of variance accounted for any illusion on any other is surprisingly small.  The second is that individual differences are related by a common factor which may even reflect personality differences \parencite[]{Makowski.etal.2023}.  We find support for both positions.  Visual illusions are undoubtedly weakly related.   Individual differences in any one illusion account for 5\% of the variance in another.   And, visual illusions are connected by a common factor that accounts for 23.3\% of the variance in each illusion when one version of each is entered into the analysis.  Of course, both proportions of variance are correct---this is how latent variables work.  One task may explain $a_1$ of the variation in the latent variable; the latent-variable may explain $a_2$ of the variation in another task, and, assuming these proportions are independent, then the first task explains $a_1\times a_2$ of the variance in the second.

In the introduction of this paper, we raised the hope that patterns of individual differences could inform about the structure of perception and cognition that underlie these illusions.  That hope remains unrealized.  There most definitely is common variation, but it seems too small to finely decompose, and, consequently, too small to understand the perceptual and cognitive factors that mediate illusions.  Moreover, because this variation is small and loads on all tasks, it may have little to do with illusions per se.  An alternative is that the factor relates to how people engage in computer tasks or how conscientious they are \parencite{Makowski.etal.2023}.  


\subsection{Future Directions}

The set of results that strikes us the most as cognitive psychologists is comprised of version correlations.  For some illusions, a small change in setup designed to reverse the direction of the illusion does exactly this.  Consider the Zöllner illusion.  In Figure~\ref{fig:illPlots}F, the top horizontal line appears slanted such that the left end is higher than the right end.  In the alternative version (not shown), the cross lines are rotated 90 degrees.  The bias reverses as expected---the top horizontal line appears slanted such that the left end is lower than the right end.  And, the correlation is high---people who show larger biases in one version show large reverse biases in the other.  This reasonable expectation does not hold for the Brentano or the Ebbinghaus Illusions.  For instance, in the Brentano illusion (Figure~\ref{fig:illPlots}A), people set the chevron too far to the left.  In the alternative version (Figure~\ref{fig:illPlots}B), they set it too far to the right.  These biases correlate moderately with other illusions, such as the Zöllner and Poggendorf.  Yet, the two versions have surprisingly little correlation between themselves.  Neither of these results is a statistical fluke as the data are highly reliable and the 95\% credible intervals on the correlation coefficients are well separated (see Figure~\ref{fig:ciCorPlots}).  How is this pattern so?  What about the Brentano and Ebbinghaus Illusions make them pathological in this regard?  Perhaps, rather than clustering illusions based on how well they correlate to each other, we should study how well they correlate within themselves, that is, how robust they are to reasonable configurable changes.  It may be this robustness that provides insight into the mechanisms after all.